% ===== 4 筛分级配映射方法 =====
\section{从视觉颗粒到筛分级配}
\label{sec:r-sieve-final}

\subsection{二维视觉粒径与筛孔粒径}

实例分割给出颗粒在相机视角下的二维投影，机械筛分则记录颗粒在振动、翻转和多姿态条件下通过给定筛孔的结果。两种尺度具有相关性，但物理定义不同。本文因此把图像几何量作为筛分粒径的观测基础，并通过显式模型定义\emph{筛分等效粒径}，使二维测量与后续级配统计之间的关系可追溯、可调整。

对第 $i$ 个颗粒，记长轴为 $a_i$、短轴为 $b_i$、投影面积为 $A_i$。面积等效直径定义为
\begin{equation}
 d_{\mathrm{eq},i}=2\sqrt{\frac{A_i}{\pi}}.
\label{eq:r-final-deq}
\end{equation}
短轴 $b_i$ 描述颗粒投影的窄向尺度，与筛孔通过具有直观几何联系；$d_{\mathrm{eq},i}$ 则把整体投影面积压缩为单一长度。细长或不规则颗粒上，两者可能产生明显差异。系统以低维线性模型同时保留这两种信息：
\begin{equation}
 d_i=\theta_1 b_i+\theta_2 d_{\mathrm{eq},i}+\theta_3,
\label{eq:r-final-sieve-d}
\end{equation}
其中 $d_i$ 为筛分等效粒径，$\boldsymbol{\theta}=(\theta_1,\theta_2,\theta_3)$ 控制二维几何到筛分统计尺度的映射。本文结果采用 $\boldsymbol{\theta}=(1,0,0)$，即 $d_i=b_i$。这一设置把短轴作为透明、可复现的基线；当前证据中没有用于拟合 $\boldsymbol{\theta}$ 的配对机械筛分数据。

\subsection{从颗粒数量到质量代理}

图像中的实例集合天然对应颗粒个数，而机械筛分通常按各筛级质量计算比例。若所有实例等权，大量小颗粒会在累计分布中占据较高权重，统计中心容易向细端移动。本文用尺寸幂律构造质量代理权重
\begin{equation}
 w_i=d_i^{\gamma},
\label{eq:r-final-mass}
\end{equation}
其中 $\gamma$ 为质量尺度指数。几何相似且材料密度近似稳定时，颗粒体积随特征长度的三次方变化，因此当前基线取 $\gamma=3$。$w_i$ 表示质量相关的统计权重，用于改变不同粒径实例在场景级分布中的贡献。

加权经验累积分布定义为
\begin{equation}
 F(d)=\frac{\sum_i w_i\,\mathbf{1}(d_i\le d)}{\sum_i w_i},
\label{eq:r-final-cdf}
\end{equation}
特征粒径由
\begin{equation}
 D_p=\inf\{d:F(d)\ge p/100\},\qquad p\in\{10,50,90\}
\label{eq:r-final-dp}
\end{equation}
得到。当前实现采用离散加权经验分布的左连续逆：累计质量代理首次达到目标比例时，对应颗粒的实际 $d_i$ 即为 $D_p$，不在相邻实例之间插值。因而任一分位数都可以回溯到具体颗粒及其权重。

\subsection{标准筛级聚合}

赛题主粒径范围为 5--40~mm。系统采用 $\{5,10,16,20,25,31.5,40\}$~mm 作为筛级边界，将每个 $d_i$ 聚合到对应区间，并保留 $<5$~mm 和 $\ge40$~mm 两个边界区间用于质量控制。第 $k$ 个粒级 $B_k$ 的质量代理占比为
\begin{equation}
 P_k=\frac{\sum_{i:d_i\in B_k}w_i}{\sum_iw_i}.
\label{eq:r-final-bin}
\end{equation}
逐颗粒几何表由此统一生成筛级比例、累计分布和 $D_{10}/D_{50}/D_{90}$。显式聚合还保留了场景级结果与颗粒级数据之间的对应关系：某一筛级比例或高分位数发生变化时，可以直接检查贡献最大的实例，而无需从黑盒回归值反推原因。

\begin{table}[t]
\centering
\caption{本文采用的筛分语义参数及其当前角色。}
\label{tab:r-final-sieve-params}
\footnotesize
\renewcommand{\arraystretch}{1.12}
\begin{tabular}{>{\raggedright\arraybackslash}p{2.3cm}>{\raggedright\arraybackslash}p{3.7cm}>{\raggedright\arraybackslash}p{6.7cm}}
\toprule
参数 & 当前设置 & 解释 \\
\midrule
$r$ & 每幅图像的 mm/px & 相机几何的物理尺度，由参考物确定 \\
$\boldsymbol{\theta}$ & $(1,0,0)$ & 以短轴作为当前筛分等效粒径基线 \\
$\gamma$ & 3 & 几何相似条件下的体积/质量尺度先验，用作质量代理权重 \\
筛级边界 & 5, 10, 16, 20, 25, 31.5, 40~mm & 固定工程统计口径，不随样本调整 \\
\bottomrule
\end{tabular}
\end{table}

\subsection{参数角色与校准接口}

$\boldsymbol{\theta}$ 决定二维几何如何压缩为筛分等效粒径，$\gamma$ 决定粗、细颗粒在场景统计中的相对权重。二者都位于实例分割之后，因此参数变化不会改变颗粒边界，只会改变物理统计口径。代码中保留了基于差分进化的校准接口：存在配对机械筛分数据时，可在 $\theta_1,\theta_2\in[0,3]$、$\theta_3\in[-20,20]$~mm 和 $\gamma\in[1,5]$ 范围内搜索，并以 $D_{10}/D_{50}/D_{90}$ 的相对误差构造目标函数。低维参数化使每个调整都能对应到粒径尺度或质量权重的明确变化，也限制了少量物理样本条件下的自由度。

本文现有三张真实图像没有配对的逐筛称量数据，因此结果统一使用短轴基线和 $d^3$ 质量代理，不执行数据驱动校准。后续实验章据此报告\emph{视觉质量代理级配}，重点分析不同场景的输出结构以及数量分布与质量代理分布之间的变化。